\section{Global Derivation Graph and Theorem Verdict}
\label{sec:derivation-graph-verdict}

The machine-readable derivation graph contains 56 nodes and represents the
seven frozen PR-I postulates as eight atomic constraint nodes.  It audits 13
major published outputs.  Parent closure, acyclicity, the diagnostic-ancestry
firewall, and the no-new-postulate freeze all pass.  The graph contains 19
blocking selectors: seven have explicit countermodels and twelve remain open.
Zero of the 13 audited outputs has a proved postulate-only ancestry.

\begin{table}[H]
\centering
\caption{Evidence-backed uniqueness status.}
\label{tab:pr35-status}
\begin{tabularx}{\textwidth}{@{}p{0.27\textwidth}X@{}}
\toprule
Block & Result \\
\midrule
PR-I radial coefficients & Unique in the declared minimal even quartic class. \\
PR-I compact closure & Not unique from P1--P7; exact cofactor and projector countermodels survive. \\
PR-II hypercharge & Unique modulo normalization and up/down relabeling only after a nonzero charge ledger is fixed. \\
PR-II photon & Unique null direction after the outer-product neutral mass matrix is fixed. \\
PR-II generations & Not unique under anomaly and Witten gates; every positive replication count survives. \\
PR-III electromagnetic tier & Not unique without a derived invariant basis and power counting. \\
PR-III $C_2$ & Not unique without a charged/neutral projector-trace theorem. \\
PR-III $D_3$ & Full real symmetric coefficient class at the fixed kernel exhausted: structural gates leave $D_3(k)=\rho kP_Z$ for every real $k$; $k=2$ is unique only under the canonical unit-prefactor physical map, whose determinant ancestry is open. \\
Provenance & No diagnostic ancestors detected in the current generative graph. \\
Global PR-I--PR-III closure & Not a singleton equivalence class under the frozen postulates. \\
\bottomrule
\end{tabularx}
\end{table}

\begin{theorem}[Frozen-postulate non-closure]
Let $\mathfrak A(P_1,\ldots,P_7)$ contain every architecture satisfying the
seven public PR-I postulates and the no-target-input methodology, modulo gauge,
basis, coordinate, phase, algebraic, and unit conventions.  Then the published
PR-I--PR-III architecture is not the only equivalence class in
$\operatorname{Sol}(\mathfrak A)$.  In particular, the exact compact cofactor
countermodels, finite odd projector countermodels, arbitrary complete
generation replications, PR-III radiative monomials, and electroweak
coefficient branches include inequivalent survivors with no empirical input
and no continuous fitted parameter.
\end{theorem}

\begin{proof}
The compact branches of Eq.~\eqref{eq:pr35-terminal-countermodels} give
distinct rational boundary factors while preserving P1--P7.  The projectors
$\Pi_{13}$ and $\Pi_{15}$ compress the nondegenerate radial operator onto
different spectral ranges and are not related by an allowed equivalence.
Complete generation replication preserves every local anomaly and Witten
parity for all positive integers.  Equation~\eqref{eq:pr35-pr3-monomial-counterexamples}
and the alternative charged/neutral integer pairs satisfy the currently stated
finite-tier dimensional and sign gates but change the generated corrections.
Thus at least two inequivalent solutions survive in several independent
sectors.  A quotient containing two inequivalent members is not a singleton.
\end{proof}

This theorem does not say that no stronger derivation could ever be found.  It
says that such a derivation is absent from, and not logically entailed by, the
seven postulates as they are currently formulated.  The freeze rule forbids
repairing this by declaring the surviving selector to be a new postulate.
